\chapter*{Prefácio}\label{prefacio}
\addcontentsline{toc}{chapter}{Prefácio}

\textbf{Uma nota de transparência}

Este guia é um instrumento analítico, não um trabalho de jornalismo investigativo nem um relatório oficial. Ele reúne dados públicos de fontes governamentais paraguaias (DGEEC, INE, INDERT, CONATEL, MEC, MSPyBS), bases internacionais (NASA POWER, PNUD, UNODC, Banco Mundial) e observações de campo sistematizadas, organizados segundo a metodologia GSS (\textit{Global Safety Score}) descrita no capítulo de Metodologia.

Toda análise contém viés. Os critérios adotados aqui privilegiam: (a) baixa densidade populacional como ativo estratégico; (b) recursos naturais renováveis e autonomia energética; (c) estabilidade jurídica e fiscal; (d) distância de alvos estratégicos em cenários de crise grave. Leitores com prioridades diferentes --- acesso a grandes mercados, vida urbana intensa, ou proximidade de centros universitários de elite --- encontrarão outras localidades mais adequadas ao seu perfil.

\medskip

\textbf{Limitações declaradas}

Os dados de segurança pública são os mais difíceis de verificar com precisão distrital. O Ministerio del Interior publica dados por departamento; a desagregação por distrito é estimada com base em relatos do Ministerio Público, cobertura jornalística e fontes secundárias. Pontuações GSS para distritos com poucos dados são conservadoras por definição.

Preços de imóveis e terra refletem cotações de 2024 e devem ser verificados antes de qualquer decisão de compra. O mercado imobiliário paraguaio tem baixa liquidez em zonas rurais e é sensível a variações cambiais do dólar americano.

\medskip

\textbf{Calibração do GSS --- Referências Mundiais}

Para dar ao leitor uma âncora concreta, a tabela abaixo posiciona o GSS de algumas localidades conhecidas internacionalmente ao lado das regiões analisadas neste guia.

\begin{center}
\small
{\renewcommand{\arraystretch}{1.3}
\begin{tabular}{p{4.5cm}cc p{3.8cm}}
\toprule
\textbf{Localidade} & \textbf{GSS} & \textbf{Faixa} & \textbf{Perfil} \\
\midrule
Idaho/Montana rural (EUA) & 9,2 & 9,0--9,5 & Remoto, autossuficiente, baixa densidade \\
Suíça rural (cantões alpinos) & 9,0 & 8,8--9,2 & Neutro, cidadão armado, estabilidade máxima \\
Nova Zelândia --- Ilha do Sul & 8,7 & 8,5--9,0 & Isolamento geográfico, recursos abundantes \\
Uruguai rural (interior) & 7,8 & 7,5--8,0 & Estabilidade institucional, baixa violência \\
\textbf{Filadelfia/Boquerón (PY)} & \textbf{7,2} & \textbf{7,0--7,5} & \textbf{Chaco Central, água escassa, muito isolado} \\
\textbf{Encarnación/Itapúa (PY)} & \textbf{6,8} & \textbf{6,5--7,0} & \textbf{Fronteira tranquila, boa infraestrutura} \\
\textbf{San Bernardino/Cordillera (PY)} & \textbf{6,7} & \textbf{6,5--7,0} & \textbf{Próximo a Asunção, área turística} \\
\textbf{Villarrica/Guairá (PY)} & \textbf{6,5} & \textbf{6,0--7,0} & \textbf{Boa infraestrutura, distância de fronteiras} \\
\textbf{Coronel Oviedo/Caaguazú (PY)} & \textbf{6,3} & \textbf{6,0--6,5} & \textbf{Eixo logístico, área de transição} \\
Bolívia oriental rural & 5,7 & 5,5--6,0 & Recursos naturais bons, instabilidade política \\
Nordeste rural do Brasil & 5,0 & 4,5--5,5 & Seca, pobreza, violência rural \\
\textbf{Concepción/Concepción (PY)} & \textbf{4,8} & \textbf{4,5--5,5} & \textbf{EPP e narcotráfico, fronteira de risco} \\
Ciudad del Este/Alto Paraná (PY) & 4,5 & 4,0--5,0 & Tríplice fronteira, criminalidade elevada \\
Caracas, Venezuela & 2,5 & 2,0--3,0 & Colapso institucional, violência extrema \\
\bottomrule
\end{tabular}}
\end{center}

\medskip
\textit{A tabela acima é um instrumento pedagógico de calibração. Os valores para localidades fora do Paraguai são estimativas baseadas nos critérios GSS aplicados a dados públicos disponíveis, e não devem ser interpretados como avaliações definitivas dessas regiões.}

\medskip

\textbf{Como este guia deve ser usado}

As pontuações GSS são pontos de partida para investigação, não vereditos finais. Nenhuma localidade com GSS 7,0 é garantia de segurança; nenhuma com GSS 5,0 é necessariamente inviável. O score resume um conjunto de variáveis estruturais. Fatores pessoais --- redes de relacionamento local, capacidade de adaptação, recursos financeiros --- podem compensar deficiências estruturais ou ampliar vantagens.

Recomenda-se que o leitor faça ao menos uma visita presencial às localidades de interesse antes de tomar qualquer decisão de compra ou mudança permanente.

\clearpage
